\Askhsh{1827}{4}
\begin{ylh}{1}
    \item 5.2. Παραλληλόγραμμα
    \item 5.7. Βαρύκεντρο τριγώνου.
\end{ylh}
Δίνεται τρίγωνο $\mathrm{A}\mathrm{B}\Gamma$ και $\mathrm{E}$ το μέσο της διαμέσου $\mathrm{B}\Delta$. Στην προέκταση της $\mathrm{A}\mathrm{E}$ θεωρούμε σημείο $\mathrm{Z}$ τέτοιο ώστε $\mathrm{E}\mathrm{Z} = \mathrm{A}\mathrm{E}$ και έστω $\Theta$ το σημείο τομής της $\mathrm{A}\mathrm{Z}$ με την πλευρά $\mathrm{B}\Gamma$.
Να αποδείξετε ότι:
\begin{alist}
    \item Το τετράπλευρο $\mathrm{A}\mathrm{B}\mathrm{Z}\Delta$ είναι παραλληλόγραμμο. \monades{8}
    \item Το τετράπλευρο $\mathrm{B}\Delta\Gamma\mathrm{Z}$ είναι παραλληλόγραμμο. \monades{8}
    \item Το σημείο $\Theta$ είναι βαρύκεντρο του τριγώνου $\mathrm{B}\Delta\mathrm{Z}$. \monades{9}
\end{alist}
